% !PRINTED 89
\weberpage{89}
%<<CONT
wählen. Dies ist eine eindeutige Funktion $F(z)$ der Variablen
\[
z = e^{2\pi iv},
\]
und wenn wie früher $q = e^{\pi i\omega}$ ist, so hat sie die Periodeneigenschaft:
\weq{2}{F(q^2z) = e^{-2\pi ia}F(z),}
woraus für jedes ganzzahlige $\nu$ folgt:
\[
F(q^{2\nu}z) = e^{-2\pi ia\nu}F(z).
\]
Die Funktion $F(z)$ wird für ein endliches $z$ nur dann unendlich,
wenn einer der Faktoren $q^{2\nu}z - 1$ verschwindet, und es
ist insbesondere
\weq{3}{(z-1)F(z) = 1 \quad (\text{für } z = 1).}
Nach (2) ist aber
\[
(q^{2\nu}z - 1)F(q^{2\nu}z) = e^{-2\pi ia\nu}F(z)(q^{2\nu}z - 1)
\]
und folglich ist nach (3)
\weq{4}{F(z)(q^{2\nu}z - 1) = e^{2\pi ia\nu} \quad (\text{für } z = q^{-2\nu}).}
Setzen wir unter der gleich noch näher zu prüfenden Voraussetzung
 der Konvergenz
\weq{5}{S(z) = \sum_{-\infty,\infty}^{\nu} \frac{e^{2\pi ia\nu}}{q^{2\nu}z - 1},}
so ergibt sich, indem wir $z$ durch $q^2z$ und $\nu$ durch $\nu - 1$ ersetzen:
\weq{6}{S(q^2z) = e^{-2\pi ia}S(z),}
und die Differenz $F(z) - S(z)$ wäre also, als Funktion von $v$
betrachtet, nach (3) und (4) eine ganze doppeltperiodische Funktion
 zweiter Art. Eine solche muß aber nach § 16, 6. eine Exponentialfunktion
 sein, und es ist also
\[
F(z) = S(z) + Ce^{-2\pi imv},
\]
worin $m$ eine ganze Zahl und $C$ eine Konstante ist. Aus (2) und
(6) aber ergibt sich, wenn nicht $C = 0$ ist:
\[
e^{-2\pi im\omega} = e^{-2\pi ia};
\]
also $a = m\omega + n$, worin $m, n$ ganze Zahlen sind. Ist aber $a$
eine Periode, so reduziert sich $F(z)$ auf eine Konstante oder eine
Exponentialfunktion. Anderenfalls muß $C = 0$ sein, und es folgt
die Entwickelung:
\weq{7}{\frac{\vartheta'_{11}\vartheta_{11}(v+a)}{2i\pi\vartheta_{11}(v)\vartheta_{11}(a)} = \sum_{-\infty,+\infty}^{\nu} \frac{e^{2\pi i\nu a}}{q^{2\nu}e^{2\pi iv} - 1}.}
